\documentclass{beamer}
\usetheme{Madrid}
\usecolortheme{default}
\usepackage[utf8]{inputenc}
\usepackage{amsmath}
\usepackage{amsfonts}
\usepackage{amssymb}
\usepackage{graphicx}
\usepackage{listings}
\usepackage{xcolor}
\usepackage{hyperref}

% Custom colors
\definecolor{codegreen}{rgb}{0,0.6,0}
\definecolor{codegray}{rgb}{0.5,0.5,0.5}
\definecolor{codepurple}{rgb}{0.58,0,0.82}
\definecolor{backcolour}{rgb}{0.95,0.95,0.92}

% Code styling
\lstdefinestyle{mystyle}{
	backgroundcolor=\color{backcolour},   
	commentstyle=\color{codegreen},
	keywordstyle=\color{magenta},
	numberstyle=\tiny\color{codegray},
	stringstyle=\color{codepurple},
	basicstyle=\ttfamily\footnotesize,
	breakatwhitespace=false,         
	breaklines=true,                 
	captionpos=b,                    
	keepspaces=true,                 
	numbers=left,                    
	numbersep=5pt,                  
	showspaces=false,                
	showstringspaces=false,
	showtabs=false,                  
	tabsize=2
}

\lstset{style=mystyle}

% Title information
\title{Mastering SQL Joins with Bridge Tables}
\subtitle{A Practical Guide for Working with Unfamiliar Data}
\author{Data Engineering Team}
\institute{Enterprise Data Solutions}
\date{\today}

\begin{document}
	
	% Slide 1: Title
	\begin{frame}
		\titlepage
	\end{frame}
	
	% Slide 2: Agenda
	\begin{frame}{Agenda}
		\begin{enumerate}
			\item Introduction to Complex Joins
			\item Understanding Bridge Tables
			\item Data Discovery Process
			\item Multi-Table Join Strategies
			\item Data Volume Handling
			\item Python \& PySpark Integration
			\item Working with Limited Information
			\item Performance Optimization
			\item Real-World Example
			\item Q\&A
		\end{enumerate}
	\end{frame}
	
	% Slide 3: Introduction
	\begin{frame}{Introduction: The Challenge}
		\begin{block}{Working with Unfamiliar Data}
			\begin{itemize}
				\item Unknown table structures and relationships
				\item Inconsistent naming conventions
				\item Limited documentation
				\item Complex business logic
				\item Joining 6-7+ tables regularly
			\end{itemize}
		\end{block}
		\vspace{0.5cm}
		\begin{alertblock}{Key Objective}
			Develop a systematic approach to navigate these challenges using modern tools like Copilot.
		\end{alertblock}
	\end{frame}
	
	% Slide 4: Bridge Tables Overview
	\begin{frame}{Understanding Bridge Tables}
		\begin{columns}
			\column{0.5\textwidth}
			\begin{block}{What is a Bridge Table?}
				\begin{itemize}
					\item Junction/Associative tables
					\item Facilitate many-to-many relationships
					\item Store relationship metadata
				\end{itemize}
			\end{block}
			\column{0.5\textwidth}
			\begin{exampleblock}{Common Example}
				\begin{itemize}
					\item Orders $\leftrightarrow$ Products
					\item Students $\leftrightarrow$ Courses
					\item Patients $\leftrightarrow$ Doctors
				\end{itemize}
			\end{exampleblock}
		\end{columns}
		\vspace{0.5cm}
		\begin{figure}
			\centering
			\begin{tabular}{|c|c|}
				\hline
				\textbf{Table} & \textbf{Purpose} \\
				\hline
				Product & Product info \\
				Order & Order header \\
				OrderProduct & Bridge table \\
				\hline
			\end{tabular}
			\caption{Bridge Table Structure}
		\end{figure}
	\end{frame}
	
	% Slide 5: Bridge Table Example
	\begin{frame}[fragile]{Bridge Table SQL Example}
		\begin{lstlisting}[language=SQL, caption=Typical Bridge Table]
			CREATE TABLE OrderProduct (
			OrderID INT,
			ProductID INT,
			Quantity INT,
			UnitPrice DECIMAL(10,2),
			PRIMARY KEY (OrderID, ProductID),
			FOREIGN KEY (OrderID) REFERENCES Orders(OrderID),
			FOREIGN KEY (ProductID) REFERENCES Products(ProductID)
			);
		\end{lstlisting}
		\begin{itemize}
			\item Composite primary key
			\item References both parent tables
			\item Contains relationship-specific attributes
		\end{itemize}
	\end{frame}
	
	% Slide 6: Data Discovery
	\begin{frame}[fragile]{Step 1: Data Discovery}
		\begin{block}{Discovery Actions}
			\begin{enumerate}
				\item Identify available tables
				\item Examine sample data
				\item Check data types
				\item Review constraints
			\end{enumerate}
		\end{block}
		\begin{lstlisting}[language=SQL, caption=Discovery Queries]
			-- For Teradata
			SELECT * FROM dbc.tables 
			WHERE databasename = 'YOUR_DB';
			
			-- Check table structure
			SHOW TABLE YOUR_TABLE;
			
			-- Sample data exploration
			SELECT * FROM YOUR_TABLE SAMPLE 100;
		\end{lstlisting}
	\end{frame}
	
	% Slide 7: Multi-Table Join Strategy
	\begin{frame}{Join Path Identification}
		\begin{alertblock}{Guidelines for 6-7 Table Joins}
			\begin{itemize}
				\item Start with fact tables
				\item Join dimension tables through bridge tables
				\item Use LEFT JOIN for optional relationships
				\item Document join logic carefully
			\end{itemize}
		\end{alertblock}
		\vspace{0.3cm}
		\begin{exampleblock}{Best Practices}
			\begin{enumerate}
				\item Use meaningful aliases
				\item Break queries with CTEs
				\item Test incrementally
				\item Validate results after each join
			\end{enumerate}
		\end{exampleblock}
	\end{frame}
	
	% Slide 8: Complex Join Example
	\begin{frame}[fragile]{Multi-Table Join Example}
		\begin{lstlisting}[language=SQL, caption=Complex Join with CTE]
			WITH order_summary AS (
			SELECT 
			o.OrderID,
			o.OrderDate,
			o.CustomerID,
			COUNT(op.ProductID) AS ProductCount
			FROM Orders o
			LEFT JOIN OrderProduct op 
			ON o.OrderID = op.OrderID
			GROUP BY o.OrderID, o.OrderDate, o.CustomerID
			)
			SELECT 
			c.CustomerName,
			os.OrderDate,
			p.ProductName,
			op.Quantity,
			p.Price * op.Quantity AS TotalAmount
			FROM order_summary os
			JOIN Customers c ON os.CustomerID = c.CustomerID
			JOIN OrderProduct op ON os.OrderID = op.OrderID
			JOIN Products p ON op.ProductID = p.ProductID
			WHERE os.ProductCount > 1;
		\end{lstlisting}
	\end{frame}
	
	% Slide 9: Data Volume Strategy
	\begin{frame}{Data Volume Strategy}
		\begin{table}
			\centering
			\begin{tabular}{|c|c|c|}
				\hline
				\textbf{Volume} & \textbf{Tool} & \textbf{Strategy} \\
				\hline
				100K-300K & Pandas & In-memory \\
				300K-2M & PySpark & Distributed \\
				>2M & Teradata & DB-side \\
				\hline
			\end{tabular}
			\caption{Processing Strategy by Data Volume}
		\end{table}
		\vspace{0.5cm}
		\begin{itemize}
			\item Consider performance vs. flexibility
			\item Balance between SQL and Python
			\item Choose based on business requirements
		\end{itemize}
	\end{frame}
	
	% Slide 10: Pandas Integration
	\begin{frame}[fragile]{Pandas Integration}
		\begin{lstlisting}[language=Python, caption=Extracting to Pandas]
			import pandas as pd
			import teradatadf
			
			conn = teradatadf.connect(
			host='your_host',
			username='your_user',
			password='your_pass'
			)
			
			query = """
			SELECT * FROM large_table
			WHERE conditions
			"""
			df = pd.read_sql(query, conn)
			
			# Further processing
			result = df.groupby('category').agg({
				'sales': 'sum',
				'quantity': 'count'
			})
		\end{lstlisting}
		\begin{itemize}
			\item Best for medium datasets
			\item Flexible analysis
			\item Rich ecosystem
		\end{itemize}
	\end{frame}
	
	% Slide 11: PySpark Integration
	\begin{frame}[fragile]{PySpark for Larger Datasets}
		\begin{lstlisting}[language=Python, caption=PySpark Integration]
			from pyspark.sql import SparkSession
			import teradatasql
			
			spark = SparkSession.builder \\
			.appName("Teradata_Integration") \\
			.getOrCreate()
			
			df_spark = spark.read \\
			.format("jdbc") \\
			.option("url", "jdbc:teradata://host") \\
			.option("dbtable", "large_table") \\
			.option("user", "username") \\
			.option("password", "password") \\
			.load()
			
			result_spark = df_spark.groupBy("category") \\
			.sum("sales") \\
			.filter("sum(sales) > 1000")
		\end{lstlisting}
	\end{frame}
	
	% Slide 12: Working with Limited Information
	\begin{frame}[fragile]{Handling Unknown Data}
		\begin{block}{Exploratory Techniques}
			\begin{itemize}
				\item Column profiling
				\item Relationship discovery
				\item Business context gathering
				\item Incremental testing
			\end{itemize}
		\end{block}
		\begin{lstlisting}[language=SQL, caption=Defensive Querying]
			-- Check for nulls
			SELECT 
			column_name,
			COUNT(*) AS total,
			SUM(CASE WHEN column_name IS NULL 
			THEN 1 ELSE 0 END) AS null_count
			FROM your_table;
		\end{lstlisting}
	\end{frame}
	
	% Slide 13: Performance Optimization
	\begin{frame}{Performance Tips}
		\begin{columns}
			\column{0.5\textwidth}
			\begin{block}{Indexing}
				\begin{itemize}
					\item Join column indexes
					\item Covering indexes
					\item Partition strategy
				\end{itemize}
			\end{block}
			\column{0.5\textwidth}
			\begin{block}{Query Optimization}
				\begin{itemize}
					\item EXPLAIN PLAN
					\item Early filtering
					\item Avoid SELECT *
					\item UNION vs UNION ALL
				\end{itemize}
			\end{block}
		\end{columns}
		\vspace{0.5cm}
		\begin{alertblock}{Memory Consideration}
			\begin{itemize}
				\item Monitor memory usage
				\item Use batch processing for large joins
				\item Consider temp tables
			\end{itemize}
		\end{alertblock}
	\end{frame}
	
	% Slide 14: Real-World Example
	\begin{frame}{Real-World Example: Sales Analysis}
		\begin{block}{Business Requirement}
			Analyze sales trends with customer demographics and product categories across 7 tables.
		\end{block}
		\begin{enumerate}
			\item Data Discovery: Identify 7 tables
			\item Join Planning: Map relationships
			\item Query Development: Build incrementally
			\item Performance Tuning: Optimize for volume
		\end{enumerate}
	\end{frame}
	
	% Slide 15: Full Query Example
	\begin{frame}[fragile]{Complete Sales Analysis Query}
		\begin{lstlisting}[language=SQL, caption=Complex Sales Analysis]
			WITH sales_data AS (
			SELECT s.SaleID, s.SaleDate, 
			s.CustomerID, s.ProductID,
			s.Quantity, s.SaleAmount,
			sb.RegionID, sb.StoreID
			FROM Sales s
			JOIN SalesBridge sb ON s.SaleID = sb.SaleID
			WHERE s.SaleDate >= DATE '2024-01-01'
			),
			customer_info AS (
			SELECT c.CustomerID, c.CustomerName,
			c.AgeGroup, c.IncomeLevel,
			r.RegionName
			FROM Customers c
			LEFT JOIN Regions r ON c.RegionID = r.RegionID
			)
			SELECT ci.RegionName, p.CategoryName,
			DATE_TRUNC('month', sd.SaleDate) AS SaleMonth,
			SUM(sd.SaleAmount) AS TotalRevenue,
			COUNT(DISTINCT sd.CustomerID) AS UniqueCustomers
			FROM sales_data sd
			JOIN customer_info ci ON sd.CustomerID = ci.CustomerID
			JOIN Products p ON sd.ProductID = p.ProductID
			JOIN Categories cat ON p.CategoryID = cat.CategoryID
			GROUP BY ci.RegionName, p.CategoryName,
			DATE_TRUNC('month', sd.SaleDate)
			ORDER BY SaleMonth DESC, TotalRevenue DESC;
		\end{lstlisting}
	\end{frame}
	
	% Slide 16: Key Takeaways
	\begin{frame}{Key Takeaways}
		\begin{enumerate}
			\item \textbf{Thorough Discovery} - Know your data before joining
			\item \textbf{Strategic Planning} - Map relationships carefully
			\item \textbf{Appropriate Tooling} - Choose based on data volume
			\item \textbf{Incremental Development} - Build and test progressively
			\item \textbf{Performance Awareness} - Optimize early and often
		\end{enumerate}
		\vspace{0.5cm}
		\begin{exampleblock}{Modern Assistant}
			Leverage tools like Copilot for assistance while maintaining understanding of underlying data relationships and business context.
		\end{exampleblock}
	\end{frame}
	
	% Slide 17: Summary
	\begin{frame}{Summary}
		\begin{block}{Complete Framework}
			\begin{itemize}
				\item Systematic approach to complex joins
				\item Bridge table mastery
				\item Volume-appropriate strategies
				\item Python integration patterns
				\item Performance optimization
			\end{itemize}
		\end{block}
		\vspace{0.3cm}
		\begin{alertblock}{Success Factors}
			\begin{itemize}
				\item Data understanding
				\item Technical flexibility
				\item Business context
				\item Continuous learning
			\end{itemize}
		\end{alertblock}
	\end{frame}
	
	% Slide 18: References
	\begin{frame}{References}
		\begin{enumerate}
			\item Teradata SQL Documentation
			\item Pandas Documentation - Data Processing
			\item PySpark Documentation - Distributed Computing
			\item SQL Performance Tuning Best Practices
			\item Database Design Patterns for Enterprise Applications
		\end{enumerate}
		\vspace{0.5cm}
		\begin{center}
			\Large{\textbf{Thank You!}} \\
			\vspace{0.3cm}
			\normalsize Questions?
		\end{center}
	\end{frame}
	
\end{document}